\subsection{Matrix-free orbital response for larger valence spaces}
For the pigment examples, the constrained ROKS orbital residual is solved
by Newton--Krylov iterations. An analytic Jacobian action differentiates the
native Fock and spin response, the orbital rotations, and the determinant
ground-overlap constraint. The covariant orbital action is symmetric; the
ground constraint is differentiated with singular-safe determinant
cofactors. No dense orbital Jacobian is allocated and no finite difference
enters the production orbital solver. All nonredundant orbital rotations
are retained. A residual-decreasing line search controls each accepted
step, while checkpointed overlap tracking follows the coherent singlet
between accepted steps. Its character relative to the initial transition
is reported separately. Final CI targets are assigned by collective
overlap with the frozen native-frontier singlet subspace; this allows
mixed frontier excitations without identifying a state by its energy.
When a target fragments over several CI eigenstates, roots are retained
until the smallest eigenvalue of their projector in the fixed target space
reaches 0.49. This keeps the original coverage threshold, protects each
target direction and is invariant to unitary rotations of the target basis.

A direct-energy initialization precedes the final KKT search when a
residual-norm minimization stagnates. The promoted open orbitals are
first held fixed while L-BFGS relaxes the remaining orbitals. Analytic
exponential-map pullbacks provide the optimization gradient. Ground
orthogonality is enforced through an augmented Lagrangian,
$E_{\rm ROKS}-\mu o+\kappa o^2/2$, with multiplier updates
$\mu\leftarrow\mu-\kappa o$. This is a numerical constraint objective;
the penalty is excluded from reported energies. All nonredundant
rotations are subsequently released and the original full-functional
KKT equations determine convergence. This freeze-and-release sequence
is motivated by the direct-optimization strategy of Levi and
co-workers.\cite{freeze2026} A second converged root is distinguished
from an existing solution by its coherent singlet overlap. If two
frontier promotions collapse onto the same root, a new initial hole
is chosen as the orthogonal complement of the converged hole within
the same frozen native frontier block. This changes the initial guess,
not the energy model or the final stationarity equations.

The large-system XC response is evaluated on its full frozen quadrature.
A density-fitted approximation preconditions the Newton linear system.
When its step fails the full-response inexact-Newton tolerance, flexible
GMRES refines it with analytic full-quadrature Hessian actions. The
ground-constraint factorization is cached within each linear solve.
The moving-frame Jacobian is also available to GMRES by subtracting the
analytic skew connection from the covariant action. The physical orbital
Hessian remains the symmetric covariant response. A Newton step that
already exceeds the trust radius need not be refined before the bounded
subspace problem; its full-response descent direction and predicted
reduction are still evaluated explicitly. For
large Newton steps, a residual-minimizing trust subspace contains the
Newton and full-response squared-residual descent directions, augmented
by directions retained from the auxiliary Krylov solve. Twice-applied
orthogonalization removes redundant directions. When retraction is enabled, every orbital basis direction is first
projected to satisfy $\mathbf j^{\mathsf T}\delta\boldsymbol\theta=0$.
The small constrained least-squares problem bounds the orbital rotation;
the unrestricted multiplier is eliminated exactly by projection in the
residual space. A thin singular-value decomposition of the response
images solves the trust problem directly. Normal equations would square
the condition number and could discard physically relevant soft modes.
The predicted reduction therefore describes the same tangent step that
the subsequent geometric retraction executes. The auxiliary response can
propose an enlarged-subspace step; its predicted reduction is then
recomputed with the full quadrature action, and acceptance requires an
actual decrease of the full functional residual. The trust radius adapts
to the ratio of actual and predicted squared-residual reduction. This
avoids excessively small steps along soft orbital modes without changing
the energy or stationarity equations. The adapted radius is saved with
the accepted orbital checkpoint and restored on continuation.
For the constrained search, each trial can additionally be retracted onto
the original zero ground-overlap manifold by short orbital rotations
along its analytic cofactor gradient. No energy evaluation enters this
geometric correction. At fixed trial orbitals the scalar multiplier is
eliminated by least squares, $\mu=(\mathbf j^{\mathsf T}\mathbf g)/
(\mathbf j^{\mathsf T}\mathbf j)$ when the constraint gradient
$\mathbf j$ is nonzero. The full KKT residual, including the overlap,
remains the convergence criterion. This avoids rejecting an otherwise
useful tangent step solely because its finite rotation creates a
second-order overlap; it does not replace ROKS by a projected energy.
Writing $a=c,s$ for charge and spin channels, the actual energy is
\begin{equation}
 E=E_{\mathrm{RI}}+\frac12\sum_a
 \left[\int \delta\bm\rho_a(\mathbf r)^{\mathsf T}
 f_{\mathrm{xc},a}(\mathbf r)\delta\bm\rho_a(\mathbf r)\,d\mathbf r
 -\mathbf q_a^{\mathsf T}K_{\mathrm{xc},a}^{\mathrm{RI}}\mathbf q_a\right],
\end{equation}
where $\bm\rho$ contains the density and its three Cartesian derivatives.
The charge deviation is measured from the anchored native ground density;
the spin deviation is the spin density itself. Thus the approximate XC
quadratic form is subtracted and the full quadrature form restored in
both the energy and the convergence residual. Stored basis values,
quadrature weights and XC kernels suffice for evaluation without a DFT
backend. On ethylene, energy, Fock and response equivalence errors are
below $2.5\times10^{-15}$ in the corresponding atomic units; the spin-pure
CI matrix agrees within $3.6\times10^{-15}$ Ha. These controls distinguish
an acceleration of the stationarity search from a change of energy model.

CI assembly likewise preserves the Hamiltonian while avoiding a full
four-index molecular-orbital transformation: the common closed-shell core
is contracted first and only the active integral block is transformed.
Fixed-bra Schur contractions reuse the common core in nonorthogonal
couplings; repeated identical canonical occupations reuse their exact
cofactors and pair contractions through a bounded cache. Retained-CSF
transition-density actions are batched into exact matrix products; identically
zero second cofactors bypass their pair-tensor contraction. Canonical columns
with a common retained bra are evaluated in batches with the same determinant
cofactors and spin coefficients. Their active/null pair features remain outer
products: the auxiliary factors are transformed along the two orbital indices
before contraction, avoiding a flattened dense feature transform. Coulomb
feature contractions use the existing auxiliary factors of the same
screened interaction, with no additional fitting or truncation. The canonical
determinant-to-CSF transformation uses its exact sparse coefficients
instead of dense multiplication; no determinant or spin coupling is
dropped. An optional Numba-compiled integer-mask Slater--Condon loop
retains the Python reference implementation; it uses neither fast-math
nor integral screening. Independent tests, including the highest bit of
a 64-spin-orbital mask, agree within $2.1\times10^{-13}$ Ha.
Singular cofactors are evaluated as products of the remaining
singular values without division. Registered density spans accelerate the
quadratic native energy only when the requested density lies in that span;
otherwise evaluation falls back to the full expression.
Frontier records are enumerated directly from at most two occupied holes
and two virtual particles, in $O(N^4)$ operations, instead of generating
and rejecting higher-rank occupations. The canonical block is retained
as dense arrays rather than a dictionary of scalar pairs. Identical
nonorthogonal matrix elements can be reused between nested windows only
when their model, screening and orbital identities match exactly.
Ground-state removal is an exact rank-one projection. Generalized
Lanczos iterations use a sparse metric inverse and its exact rank-one
Sherman--Morrison update; only the requested roots are returned, without
reducing the configuration space. On the archived 2,145-CSF problem,
the lowest twenty energies reproduce the full dense solution within
$2.0\times10^{-12}$ eV, with residuals below $1.5\times10^{-13}$ Ha.
Observable contractions use exactly the determinant ordering of the
Hamiltonian. Phase tests cover even and odd occupations and direct
overlaps with the stored metric. Identical native diagonal and four-spin
functional values can be reused when only the screened interaction
changes; an independent matrix reconstruction agrees within
$2.0\times10^{-11}$ Ha for the archived carotene block and
$1.1\times10^{-10}$ Ha for chlorophyllide. The 45 retained-CSF scalar
elements of chlorophyllide reproduce the unbatched journal within
$1.9\times10^{-12}$ Ha. Batched reconstruction of the 3,406-CSF
chlorophyllide pencil at $N=W=9$ agrees with the saved matrix within
$6.9\times10^{-13}$ Ha. Positive definiteness is certified by a Gershgorin lower bound or by
eliminating the sparse canonical block and testing its small border
Schur complement. A smallest-eigenvalue calculation remains the fallback.


Native real d functions are included in the molecular-frame transport
with analytic first and second rotation derivatives. Large-system native
response also uses the same D4 and electrostatic cutoffs as its ground
functional. Dense pair tensors and screened interactions are still stored:
matrix-free orbital stationarity does not make this implementation
linear-scaling or inexpensive in memory. The present quadrature adapter
supports vertical energies and orbital response only; it explicitly rejects
nuclear-derivative requests. The analytic nuclear Hessian qualifications
reported for the small-molecule anchored surface therefore do not extend
to these pigment calculations.
